\documentclass[10pt,aspectratio=169]{beamer}

% --- Minimal theme ---------------------------------------------------------
\usetheme{default}
\usecolortheme{dove}
\usefonttheme{professionalfonts}
\setbeamertemplate{navigation symbols}{}
\setbeamertemplate{footline}[frame number]
\setbeamertemplate{itemize items}[circle]
\setbeamertemplate{section in toc}[sections numbered]
\setbeamerfont{frametitle}{size=\large,series=\bfseries}
\setbeamercolor{frametitle}{fg=black}
\setbeamercolor{title}{fg=black}
\setbeamercolor{structure}{fg=black}

\usepackage{booktabs}
\usepackage{array}
\usepackage{graphicx}
\graphicspath{{../figs/}}
\usepackage{amsmath}

\title{Synthesising five reanalyses}
\subtitle{A rigorous comparison of ensemble-mean, median, PC1, factor model, \\
  inverse-variance, and UDel-anchored combinations}
\author{}
\date{\today}

\begin{document}

\begin{frame}\titlepage\end{frame}

\begin{frame}{Outline}\tableofcontents\end{frame}

% ===========================================================================
\section{Setup}

\begin{frame}{The problem}
  \begin{itemize}
    \item We observe five noisy measurements $\{T_k(c,t)\}_{k=1}^5$ of a
          latent true temperature $T^*(c,t)$ per (country, year).
    \item Want a synthesis $\hat T(c,t)$ with a better signal-to-noise
          ratio than any single $T_k$ --- without access to ground truth
          and without assuming errors are independent across datasets
          (ERA5 and JRA-3Q share assimilation; GMFD and UDel share station
          inputs).
    \item Throughout, the growth regression is:
      \[
      \texttt{growth} = \beta_1 T + \beta_2 T^2 + \boldsymbol\gamma^\top \mathbf{P}
          + \alpha_i + \tau_t + \delta_i t + \phi_i t^2 + u
      \]
      with clustered SE on country. Precipitation controls are
      \emph{ensemble-mean} precipitation, held fixed across methods, so the
      only thing varying is the temperature measurement.
  \end{itemize}
\end{frame}

\begin{frame}{Methods evaluated}
  Eleven temperatures enter the comparison: the five raw datasets and six
  synthetic constructions.
  \vspace{0.4em}

  \small
  \begin{tabular}{lll}
    \toprule
    Method & Construction & Assumption \\
    \midrule
    MEAN      & $\tfrac{1}{5}\sum_k T_k$                       & iid errors \\
    MEDIAN    & $\text{median}_k T_k$                          & robust (nonparametric) \\
    PC1       & first PC of residualised $T$, added to FE fit  & one-factor with constant loadings \\
    FACTOR    & per-country factor model: $R_k(c,t)=\lambda_k(c)F(c,t)+\varepsilon_k$ & factor structure, loadings vary by $c$ \\
    MIN-VAR   & $w_k(c)\propto 1/\sigma_k^2(c)$ from factor model & efficient GLS under factor model \\
    ANCHORED  & UDel as proxy truth; $w_k(c)\propto 1/\text{MSE}_k(c)$ & UDel $\approx$ observations \\
    \bottomrule
  \end{tabular}

  \vspace{0.4em}
  \footnotesize
  A seventh candidate, OOS-OPT (convex weights tuned to CV-MSE), was
  attempted and rejected --- see later.
\end{frame}

% ===========================================================================
\section{Construction details}

\begin{frame}{MEAN and MEDIAN}
  \begin{itemize}
    \item $\hat T_{\text{MEAN}} = \tfrac{1}{5}\sum_k T_k$. MVUE if errors are iid
          Gaussian. Not iid here --- errors are correlated, so effective
          variance reduction is below 1/5.
    \item MEDIAN is robust to single-dataset shocks (e.g.\ the GMFD
          2009--2010 anomaly). Asymptotically $(\pi/2)/n$ relative to MEAN
          under Gaussianity --- less efficient, but honest when
          distributional assumptions are violated.
  \end{itemize}
\end{frame}

\begin{frame}{PC1: the one-factor model, pooled}
  Let $\tilde T_k$ be the residualised temperature (after
  $\alpha_i + \tau_t + \delta_i t + \phi_i t^2$). Pool across all
  (iso, year) and run PCA on the $N \times 5$ matrix:
  \[
  \tilde T_k = \lambda_k F + \varepsilon_k, \qquad \varepsilon_k \perp F.
  \]
  First PC carries \textbf{88.0\%} of pooled residual variance. Rescale
  $\hat F$ to match the sd of the ensemble-mean residual, and add back the
  fitted FE structure to recover a °C-scale synthetic.
\end{frame}

\begin{frame}{FACTOR: per-country EM}
  Same model but loadings and idiosyncratic variances vary by country:
  \[
  \tilde T_k(c, t) = \lambda_k(c)\,F(c, t) + \varepsilon_k(c, t),
  \qquad \varepsilon_k(c,\cdot) \sim \mathcal{N}(0, \sigma_k^2(c)).
  \]
  \vspace{-0.5em}
  \begin{itemize}
    \item EM, 60 max iters, converges in $\leq 20$. Normalisation:
          $\text{sd}(F) = 1$, sign chosen so mean($\lambda$) $> 0$.
    \item Yields, as a by-product, \textbf{measurement-error variance
          estimates} $\hat\sigma_k^2(c)$ for every (dataset, country).
  \end{itemize}
\end{frame}

\begin{frame}{MIN-VAR and ANCHORED: inverse-variance combinations}
  Under the factor model, the \emph{efficient} unbiased combination is
  \[
    \hat T_{\text{MIN-VAR}}(c, t)
    = \sum_k w_k(c)\, T_k(c, t), \qquad
    w_k(c) = \frac{1/\sigma_k^2(c)}{\sum_{k'} 1/\sigma_{k'}^2(c)}.
  \]
  ANCHORED is the same idea but uses UDel as a station proxy:
  \[
  w_k(c) \propto \frac{1}{\text{MSE}\bigl(T_k(c,\cdot),\, T_{\text{UDel}}(c,\cdot)\bigr)},
  \qquad k \in \{\text{ERA5},\text{GMFD},\text{JRA-3Q},\text{MERRA2}\}.
  \]
\end{frame}

\begin{frame}{OOS-OPT: why we dropped it}
  A natural candidate: weights $w$ that minimise CV-MSE of the growth regression.
  \vspace{0.3em}

  We drew 500 weight vectors uniformly on the 5-simplex (Dirichlet(1,\ldots,1))
  and evaluated CV-MSE at each:

  \vspace{0.3em}
  \centering
  \small
  \begin{tabular}{l r}
    \toprule
    min CV-MSE (random 10-fold) & 0.002714 \\
    max CV-MSE                  & 0.002716 \\
    relative spread             & \textbf{0.050\%} of mean \\
    \bottomrule
  \end{tabular}

  \vspace{0.6em}
  \footnotesize
  The objective is essentially \emph{flat} over the simplex.  Any weights
  give (statistically) the same predictive fit, so the argmin is not
  identified. Honest conclusion: OOS-MSE cannot be used to select synthesis
  weights on this data.
\end{frame}

% ===========================================================================
\section{Evaluation framework}

\begin{frame}{Three precision criteria, two stability criteria}
  For each of 11 methods (5 raw + 6 synthetic) we report:

  \vspace{0.4em}
  \small
  \begin{itemize}
    \item \textbf{$\hat\beta_2$ magnitude} --- deattenuation check: removing
          measurement error should make the quadratic coefficient more
          negative.
    \item \textbf{CI width on ME(25°C)} --- 95\% delta-method CI at the hot
          tail, the quantity usually reported.
    \item \textbf{t-statistic on $\hat\beta_2$} --- $|\hat\beta_2/\text{SE}|$,
          the econometric measure of precision (signal-to-noise of the
          coefficient).
    \item \textbf{OOS MSE} on three CV splits (random 10-fold, leave-country,
          leave-year) --- predictive fit.
    \item \textbf{Jackknife stability}: rebuild each synthetic dropping one
          dataset, measure the range of ME(25°C).
  \end{itemize}
\end{frame}

% ===========================================================================
\section{Results}

\begin{frame}{ME curves, raw vs synthetic}
  \centering
  \includegraphics[width=\linewidth]{sc_me_curves.pdf}

  \vspace{0.2em}
  \footnotesize
  Ribbons = 95\% delta-method CIs, clustered SE on country. Synthetics
  cluster tightly; the raw panel shows more fanning, especially in the hot
  tail.
\end{frame}

\begin{frame}{Headline table (ME at 25°C)}
  \centering
  \small
  \begin{tabular}{l r r r r r}
    \toprule
    Method & $\hat\beta_2 \times 10^4$ & SE($\hat\beta_2$) $\times 10^4$
           & $|t|$ & ME(25) & CI width \\
    \midrule
    \textit{Synthetic:} & & & & & \\
    FACTOR      & $-6.28$ & 1.62 & \textbf{3.86} & $-0.0151$ & 0.0171 \\
    PC1         & $-6.40$ & 1.67 & 3.83 & $-0.0149$ & 0.0170 \\
    MEAN        & $-6.40$ & 1.68 & 3.82 & $-0.0149$ & 0.0170 \\
    ANCHORED    & $-6.04$ & 1.59 & 3.80 & $-0.0142$ & 0.0161 \\
    MEDIAN      & $-5.80$ & 1.66 & 3.50 & $-0.0131$ & 0.0166 \\
    MIN-VAR     & $-4.95$ & 1.48 & 3.34 & $-0.0107$ & 0.0158 \\
    \midrule
    \textit{Raw:} & & & & & \\
    JRA-3Q      & $-6.22$ & 1.63 & 3.81 & $-0.0161$ & 0.0165 \\
    ERA5-025    & $-5.51$ & 1.46 & 3.79 & $-0.0127$ & 0.0147 \\
    UDel        & $-5.17$ & 1.38 & 3.75 & $-0.0120$ & 0.0158 \\
    GMFD        & $-6.43$ & 1.82 & 3.53 & $-0.0112$ & 0.0162 \\
    MERRA2      & $-4.18$ & 1.41 & 2.97 & $-0.0073$ & 0.0138 \\
    \bottomrule
  \end{tabular}
\end{frame}

\begin{frame}{Coefficient plot: $\hat\beta_2$ with 95\% CI}
  \centering
  \includegraphics[height=0.75\textheight]{sc_coef_beta2.pdf}

  \vspace{0.2em}
  \footnotesize
  Sorted descending in $\hat\beta_2$.  Synthetics (blue) cluster near
  $-6\!\times\!10^{-4}$; raw datasets (grey) span from
  $-4.2\!\times\!10^{-4}$ (MERRA2, attenuated) to $-6.4\!\times\!10^{-4}$ (GMFD).
\end{frame}

\begin{frame}{Coefficient plot: ME at 25°C with 95\% CI}
  \centering
  \includegraphics[height=0.75\textheight]{sc_coef_me25.pdf}

  \vspace{0.2em}
  \footnotesize
  Same view on the ME at 25°C.  All synthetics and four of five raws reject
  zero.  MERRA2 barely does.
\end{frame}

\begin{frame}{CI width: synthetics are \emph{wider}, but that's misleading}
  \centering
  \includegraphics[height=0.55\textheight]{sc_ci_width.pdf}

  \vspace{0.2em}
  \footnotesize
  MERRA2 has the tightest CI on ME(25°C) --- but its $|\hat\beta_2|$ is
  also the smallest ($4.2 \times 10^{-4}$ vs $6.4 \times 10^{-4}$ for MEAN).
  Attenuation shrinks \emph{both} the coefficient \emph{and} its SE. The
  right econometric comparison is the t-stat on $\hat\beta_2$.
\end{frame}

\begin{frame}{t-statistic: synthetics win on the right metric}
  \small
  Ranked by $|t_{\hat\beta_2}|$:

  \vspace{0.4em}
  \centering
  \begin{tabular}{l r}
    \toprule
    FACTOR      & \textbf{3.86} \\
    PC1         & 3.83 \\
    MEAN        & 3.82 \\
    JRA-3Q      & 3.81 \\
    ANCHORED    & 3.80 \\
    ERA5-025    & 3.79 \\
    UDel        & 3.75 \\
    GMFD        & 3.53 \\
    MEDIAN      & 3.50 \\
    MIN-VAR     & 3.34 \\
    MERRA2      & 2.97 \\
    \bottomrule
  \end{tabular}

  \vspace{0.4em}
  \footnotesize
  Top three are all synthetics (FACTOR, PC1, MEAN). Gain over best raw
  is modest ($3.86$ vs $3.81$), gain over worst raw is large
  ($3.86$ vs $2.97$).
\end{frame}

\begin{frame}{Why deattenuation works, quantitatively}
  Classical measurement error in a regressor:
  \[
    \hat\beta \xrightarrow{p} \beta \cdot
      \frac{\mathrm{Var}(T^*)}{\mathrm{Var}(T^*) + \mathrm{Var}(\eta)}
  \]
  Averaging $K$ independent measurements cuts $\mathrm{Var}(\eta)$ by $1/K$.
  Errors aren't independent here, but residualised correlations 0.74--0.96
  leave plenty of idiosyncratic component to average away.

  \vspace{0.4em}

  Empirically: raw MERRA2 $\hat\beta_2 = -4.2\!\times\!10^{-4}$ vs MEAN
  $\hat\beta_2 = -6.4\!\times\!10^{-4}$. MEAN is \emph{50\% larger in
  magnitude} --- exactly the pattern deattenuation predicts.
\end{frame}

\begin{frame}{OOS predictive fit is flat --- as expected}
  \centering
  \includegraphics[width=\linewidth]{sc_oos_mse.pdf}

  \vspace{0.2em}
  \footnotesize
  All 11 methods give CV-MSE within 0.3\% of each other on all three splits.
  Leave-country-out MSE is actually slightly \emph{lower} than random
  10-fold (an artefact of the FWL leakage, symmetric across methods).
  Bottom line: predictive fit cannot discriminate among these temperatures.
\end{frame}

\begin{frame}{Jackknife stability: leave-one-dataset-out}
  \centering
  \includegraphics[width=\linewidth]{sc_jackknife.pdf}

  \vspace{0.2em}
  \small
  Range of ME(25°C) across the 5 jackknife folds:

  \vspace{0.2em}
  \centering
  \small
  \begin{tabular}{l r}
    \toprule
    FACTOR   & 0.00150 \\
    MEAN     & 0.00164 \\
    PC1      & 0.00170 \\
    ANCHORED & 0.00197 \\
    MEDIAN   & 0.00237 \\
    MIN-VAR  & 0.00321 \\
    \bottomrule
  \end{tabular}

  \vspace{0.2em}
  \footnotesize
  FACTOR is the most stable synthesis (smallest range); MIN-VAR is the
  least, because inverse-variance weights shift sharply when one dataset
  is removed.
\end{frame}

\begin{frame}{Factor-model measurement-error variances}
  \centering
  \includegraphics[width=0.82\linewidth]{sc_factor_sigma.pdf}

  \vspace{0.2em}
  \footnotesize
  $\hat\sigma_k$ per dataset, pooled across 158 country-level factor fits.
  ERA5 and JRA-3Q are nearly tied for cleanest weather signal
  (median $\approx 0.06$°C); GMFD is the noisiest ($\approx 0.20$°C).
  This is a model-based formalisation of the residual-correlation pattern
  in the previous deck.
\end{frame}

% ===========================================================================
\section{Synthesis behaviour in disagreeing countries}

\begin{frame}{Where the synthesis has to work hardest}
  Rank 158 countries by ensemble residualised MAD across all 10 pairs; take
  the top 8. These are the countries where the reanalyses disagree most on
  weather:
  \vspace{0.4em}

  \centering HTI, BDI, RWA, GRL, ARM, IRQ, GTM, ERI
\end{frame}

\begin{frame}{Raw temperatures (grey) + synthetics (colour), top-8 disagreement}
  \centering
  \includegraphics[height=0.85\textheight]{sc_spaghetti_raw.pdf}
\end{frame}

\begin{frame}{Same in residualised space (climate + trend partialled out)}
  \centering
  \includegraphics[height=0.85\textheight]{sc_spaghetti_resid.pdf}
\end{frame}

\begin{frame}{What the spaghetti plots show}
  \begin{itemize}
    \item In these most-disagreeing countries, the raw series fan out by
          0.5--1.5\,°C in individual years (grey lines).
    \item All synthetics track the central tendency of the fan, but they
          \emph{differ} among themselves in the residualised space ---
          especially FACTOR and MIN-VAR, which down-weight noisier datasets
          and therefore follow the cleaner subset.
    \item MEAN and PC1 are visually nearly identical and sit in the middle
          of the fan, as their construction implies.
    \item UDel-anchored pulls toward UDel in countries where UDel has good
          station coverage, and away from it where it doesn't.
  \end{itemize}
\end{frame}

% ===========================================================================
\section{Recommendation}

\begin{frame}{Recommendation}
  \begin{enumerate}
    \item \textbf{For the main specification: FACTOR.}
          Best $t$-statistic on $\hat\beta_2$; lowest jackknife range;
          yields measurement-error variance estimates as a by-product.
    \item \textbf{For robustness: MEAN.}
          0.1\% worse on t-stat, 9\% higher jackknife range than FACTOR,
          but indistinguishable in practice and zero tuning.
    \item \textbf{Report uncertainty bands across synthesis methods.}
          The envelope of ME curves across synthetics is narrower than the
          envelope across raw datasets --- a useful honest-uncertainty band
          for the paper.
    \item \textbf{Do \emph{not} use OOS MSE to select methods.} The
          objective is flat.  Report OOS MSE as a sanity check on
          predictive equivalence instead.
    \item \textbf{Do \emph{not} over-interpret MIN-VAR.} Its weights are
          noisy (per-country factor-model fit on 31 years), and its
          jackknife instability shows it is not robust.
  \end{enumerate}
\end{frame}

\begin{frame}{Caveats}
  \begin{itemize}
    \item \textbf{Correlated errors.} Our factor model assumes
          $\mathrm{Cov}(\varepsilon_j, \varepsilon_k | F) = 0$.  With 5 datasets
          and 31 years per country, we cannot reject this, but ERA5 and
          JRA-3Q share ECMWF-style assimilation so this is strong.
    \item \textbf{Residualisation leakage.} All methods use
          full-sample residualisation for diagnostics. Symmetric across
          methods, so rankings are unaffected; absolute $R^2_{\text{oos}}$
          values are slightly optimistic.
    \item \textbf{FWL under measurement error.} The classical deattenuation
          formula assumes classical ME in a linear regression. Our
          regression is quadratic, so the translation is approximate;
          still, the observed $\hat\beta_2$ shift is in the direction
          deattenuation predicts.
    \item \textbf{Station benchmark.} ANCHORED uses UDel as a proxy for
          ``station truth''. A proper station-anchored synthesis would use
          Berkeley Earth or GHCN-Monthly; left as a clear follow-up.
  \end{itemize}
\end{frame}

\end{document}
